\chapter{Mathematical and Cryptographic Background}
\label{ch:background}

This chapter introduces the notation and mathematical foundations
used throughout the thesis.


% ==================================================================
\section{Finite Fields}
\label{sec:finite_fields}
% ==================================================================

\begin{definition}[Prime Field]
\label{def:prime_field}
Let $p$ be a prime number. The \emph{prime field} $\Fp$ is the set
$\{0, 1, \ldots, p-1\}$ equipped with addition and multiplication
modulo $p$.
\end{definition}

The multiplicative group $\Fp^* = \Fp \setminus \{0\}$ is cyclic of
order $p - 1$. An element $g \in \Fp^*$ is called a \emph{generator}
if $\{g^0, g^1, \ldots, g^{p-2}\} = \Fp^*$. For any $a \in \Fp^*$,
the multiplicative inverse can be computed via Fermat's little theorem:
%
\begin{equation}
\label{eq:fermat}
a^{-1} \equiv a^{p-2} \pmod{p}.
\end{equation}

\begin{example}
\label{ex:prime_field}
In $\F_7 = \{0,1,\ldots,6\}$ arithmetic wraps around modulo $7$: for
instance $5 + 4 = 2$ and $5 \cdot 4 = 6$. The inverse of $3$ is
$3^{p-2} = 3^{5} = 5$, since $3 \cdot 5 = 15 = 1 \bmod 7$.
\end{example}


\subsection{The BabyBear Field}
\label{subsec:babybear}

The implementation described in this thesis operates over the BabyBear
prime field, defined by
%
\begin{equation}
\label{eq:babybear_prime}
p = 2^{31} - 2^{27} + 1 = 2\,013\,265\,921.
\end{equation}

% p < 2^31 so elements fit in 32 bits; products fit in 64 bits.
Since $p < 2^{31}$, every field element fits in a 32-bit unsigned
integer, and the product of two elements fits in a 64-bit integer.

% p - 1 = 2^{27} * 15, giving a multiplicative subgroup of order 2^{27}.
The multiplicative group $\Fp^*$ has order $p - 1 = 2^{27} \cdot 15$.
The factor $2^{27}$, the \emph{two-adicity} of the field, means that
$\Fp^*$ contains a cyclic subgroup of order $2^{27}$, which is
essential for the Number Theoretic Transform.


\subsection{Roots of Unity}
\label{subsec:roots_of_unity}

\begin{definition}[Primitive Root of Unity]
\label{def:root_of_unity}
An element $\omega \in \Fp$ is a \emph{primitive $N$-th root of unity}
if $\omega^N = 1$ and $\omega^k \neq 1$ for all $0 < k < N$.
\end{definition}

\begin{example}
\label{ex:root_of_unity}
In $\F_7$ the element $\omega = 2$ is a primitive $3$rd root of unity:
$2^1 = 2$, $2^2 = 4$, $2^3 = 1$, and no smaller power gives $1$. Its
powers form the subgroup $\{1, 2, 4\}$ of order $3$.
\end{example}

Given a primitive $N$-th root of unity $\omega$, the set
$H = \{1, \omega, \omega^2, \ldots, \omega^{N-1}\}$
is a multiplicative subgroup of $\Fp^*$ of order $N$. Because BabyBear
has two-adicity $27$, subgroups of order $2^k$ for any $k \leq 27$
exist and are generated from a fixed multiplicative generator $g$ by
taking $g^{(p-1)/2^k}$.

% ethSTARK Section 3 (Definition 1, AIR): "the following definition
% restricts our attention only to finite fields F that contain a large
% multiplicative subgroup of size 2^h".
These subgroups serve as evaluation domains for the polynomials in the
Scalable Transparent Argument of Knowledge (STARK) protocol: the execution trace is evaluated on a subgroup $H$ of
order $2^h$, and the Fast Reed-Solomon Interactive Oracle Proof of
Proximity (FRI) operates on a larger domain.


\subsection{Field Extensions}
\label{subsec:field_ext}

\begin{definition}[Field Extension]
\label{def:field_ext}
A \emph{field extension} $\mathbb{F}_{p^k}$ of degree $k$ over $\Fp$
is the quotient ring $\Fp[X]/(f(X))$, where $f(X) \in \Fp[X]$ is an
irreducible polynomial of degree $k$. Elements are polynomials of
degree less than $k$ with coefficients in $\Fp$.
\end{definition}

In this work we use the \emph{quartic extension} $\Fpk$, constructed
via the irreducible polynomial $X^4 - 11$. An element of $\Fpk$ is a
tuple $(c_0, c_1, c_2, c_3) \in \Fp^4$ representing
$c_0 + c_1 X + c_2 X^2 + c_3 X^3$, with the reduction rule $X^4 = 11$.
The extension has $p^4 \approx 2^{124}$ elements, large enough that a
folding challenge sampled from it makes the soundness error negligible,
as discussed in \cref{ch:fri}.

\begin{example}
\label{ex:field_ext}
A smaller instance is $\F_{49} = \F_7[X]/(X^2 + 1)$, which is a field
because $X^2 + 1$ has no root in $\F_7$. Its elements are $a + bX$ with
$a, b \in \F_7$ and the rule $X^2 = -1$. For example,
$(2 + 3X)(1 + X) = 2 + 5X + 3X^2 = 2 + 5X - 3 = 6 + 5X$.
\end{example}


% ==================================================================
\section{Polynomials}
\label{sec:polynomials}
% ==================================================================

A univariate polynomial over $\Fp$ of degree at most $d$ is
$f(X) = \sum_{i=0}^{d} c_i X^i$ with $c_i \in \Fp$. It can be
represented by its \emph{coefficient vector} $(c_0, \ldots, c_d)$ or
by its \emph{evaluation vector} on a domain $D$. Converting between
the two representations, and evaluating a polynomial at a point, are
the fundamental operations of the prover.


\subsection{The Number Theoretic Transform}
\label{subsec:ntt}

The conversion between coefficient and evaluation representations is
performed using the Number Theoretic Transform (NTT), the finite-field
analogue of the Fast Fourier Transform. When the evaluation domain is
a multiplicative subgroup
$D = \{1, \omega, \ldots, \omega^{N-1}\}$ with $N$ a power of~$2$, the
NTT computes all evaluations $f(\omega^i)$ in $O(N \log N)$ operations.
The algorithm exploits the decomposition
%
\[
f(X) = f_{\mathrm{even}}(X^2) + X \cdot f_{\mathrm{odd}}(X^2),
\]
%
recursing on the even- and odd-indexed coefficients over the squared
subgroup, which has half the size. The inverse transform runs the same
routine with $\omega^{-1}$ and divides by $N$.


\subsection{Low-Degree Extension}
\label{subsec:lde}

The \emph{low-degree extension} (LDE) of a polynomial is its
evaluation on a domain larger than its coefficient count. Given the
$n$ coefficients of $f$, the LDE with blowup factor $b$ evaluates $f$
on a coset of a subgroup of size $b \cdot n$. The resulting evaluation
vector is the Reed-Solomon codeword that the prover commits to. In our
implementation the low-degree extension is computed by padding the
coefficients with zeros, applying a coset shift, and running the NTT
on the enlarged domain~\cite[Section~3.4]{ethSTARK}.

\subsection{Vanishing Polynomials}
\label{subsec:vanishing}

% ethSTARK, before Definition 1 (AIR): "Given a set S in F, we define
% the vanishing polynomial of S to be Z_S(X) := prod_{s in S}(X - s)".
\begin{definition}[Vanishing Polynomial {\cite[Section~5.1]{ethSTARK}}]
\label{def:vanishing}
Given a set $S \subseteq \F$, the \emph{vanishing polynomial} of $S$
is $Z_S(X) := \prod_{s \in S}(X - s)$. When $S$ is a multiplicative
subgroup $H = \{1, \omega, \ldots, \omega^{N-1}\}$, this simplifies
to $Z_H(X) = X^N - 1$.
\end{definition}

\begin{example}
\label{ex:vanishing}
Take the subgroup $H = \{1, 2, 4\}$ of $\F_7^*$ from
\cref{ex:root_of_unity}. Its vanishing polynomial is
$Z_H(X) = (X-1)(X-2)(X-4) = X^3 - 1$, which evaluates to $0$ at each of
$1, 2, 4$ and is nonzero elsewhere, for instance $Z_H(3) = 5$.
\end{example}

$Z_H$ vanishes on every element of $H$ and nowhere else. If a
constraint polynomial $c(X)$ must be zero on $H$, then $Z_H(X)$
divides $c(X)$, and the quotient $q(X) = c(X)/Z_H(X)$ is a
polynomial. The verifier checks that $q$ has low degree, which
certifies that $c$ vanishes on $H$~\cite{ethSTARK}. This mechanism is
the basis of the STARK construction.

% ==================================================================
\section{Coding Theory}
\label{sec:coding_theory}
% ==================================================================

We recall the standard notions from coding theory that the soundness
analysis of the FRI protocol relies on. The specific family of codes
used in this work, the Reed-Solomon codes, is introduced in
\cref{sec:reed_solomon} as an instance of these general definitions.

\begin{definition}[Error-Correcting Code {\cite[Definition~3.1]{WHIR24}}]
\label{def:code}
An \emph{error-correcting code} of length $n$ over an alphabet
$\Sigma$ is a subset $C \subseteq \Sigma^n$. The code is a
\emph{linear code} if $\Sigma = \F$ is a field and $C$ is a linear
subspace of $\F^n$. The \emph{rate} of a linear code is
$\rho := \dim(C)/n$.
\end{definition}

\begin{example}
\label{ex:code}
The set of all vectors $(a, b, a+b) \in \F_7^3$ is a linear code of
length $n = 3$: it is a subspace of dimension $2$ (the free choices $a$
and $b$), so its rate is $\rho = 2/3$. The third coordinate is a
built-in check that lets some errors be detected.
\end{example}

The elements of $C$ are its \emph{codewords}, and the \emph{minimum
distance} of a code is the smallest relative Hamming distance
(\cref{def:hamming}) between two distinct codewords. When a word lies
far from every codeword, several codewords may still lie within a given
distance of it; the following notion quantifies this.

\begin{definition}[List Decoding {\cite[Definition~3.2]{WHIR24}}]
\label{def:list_decoding}
For a code $C \subseteq \Sigma^n$, a word $f \in \Sigma^n$, and a
proximity parameter $\delta \in [0,1]$, the \emph{list} of $C$ at $f$
within radius $\delta$ is
%
\[
\Lambda(C, f, \delta) := \{u \in C : \Delta(f, u) \leq \delta\}.
\]
%
The code is \emph{$(\delta, \ell)$-list decodable} if
$|\Lambda(C, f, \delta)| \leq \ell$ for every $f \in \Sigma^n$. A
radius $\delta$ is \emph{within unique decoding} if the code is
$(\delta, 1)$-list decodable.
\end{definition}

\begin{example}
\label{ex:list_decoding}
Within the unique decoding radius there is at most one codeword near a
received word, so the list $\Lambda(C, f, \delta)$ has size at most $1$
and decoding is unambiguous. Past that radius the list can grow: a word
sitting exactly between two codewords has both of them in its list, so
the code is no longer $(\delta,1)$-list decodable there.
\end{example}

For a code of minimum distance $\delta_C$, any radius
$\delta \leq \delta_C/2$ is within unique decoding. Beyond it, the list
size is controlled, for the codes we use, by the Johnson bound stated
below.

\begin{theorem}[Johnson Bound {\cite[Theorem~4.3]{WHIR24}}]
\label{thm:johnson}
The Reed-Solomon code of rate $\rho$ is
$\bigl(1 - \sqrt{\rho} - \eta,\; \tfrac{1}{2\eta\sqrt{\rho}}\bigr)$-list
decodable for every $\eta \in (0, 1 - \sqrt{\rho})$.
\end{theorem}

\begin{example}
\label{ex:johnson}
For rate $\rho = 1/2$, the unique decoding radius is
$\tfrac{1-\rho}{2} = 0.25$, while the Johnson radius is
$1 - \sqrt{1/2} \approx 0.29$. Between these two values a received word
may have a few codewords in its list, but the Johnson bound guarantees
the number stays small, which is what the soundness analysis needs.
\end{example}

The two thresholds, the \emph{unique decoding} radius
$\tfrac{1-\rho}{2}$ and the \emph{Johnson} radius $1 - \sqrt{\rho}$,
are exactly where the soundness guarantees of FRI change form
(\cref{ch:fri}).

% ==================================================================
\section{Reed-Solomon Codes}
\label{sec:reed_solomon}
% ==================================================================

% BS+18, Section 2.1: "For S in F and rate parameter rho in (0,1), the
% Reed-Solomon code RS[F, S, rho] is the family of functions f : S -> F
% that are evaluations of polynomials of degree < rho|S|."
\begin{definition}[Reed-Solomon Code {\cite[Section~2.1]{BBHR18}}]
\label{def:rs_code}
For a finite field $\F$, a subset $S \subseteq \F$, and a rate
parameter $\rho \in (0,1)$, the \emph{Reed-Solomon code}
$\RS[\F, S, \rho]$ is the set of functions $f : S \to \F$ that are
evaluations of polynomials of degree less than $\rho|S|$:
%
\[
\RS[\F, S, \rho] = \{(f(x))_{x \in S} : f \in \F[X],\;
  \deg(f) < \rho|S|\}.
\]
\end{definition}

\begin{example}
\label{ex:rs_code}
Take $\F_7$, evaluation set $S = \{1,2,3,4,5,6\}$, and rate
$\rho = 1/3$, so codewords are evaluations of polynomials of degree
less than $\rho|S| = 2$, that is, lines $a + bX$. The polynomial
$f(X) = 2 + 3X$ gives the codeword
$(f(1),\ldots,f(6)) = (5, 1, 4, 0, 3, 6)$. Any two distinct lines agree
in at most one point, so their codewords differ in at least five of the
six positions.
\end{example}

Each codeword is the evaluation of a polynomial of degree less than
$\rho|S|$ on all points of $S$. The code has minimum distance
$|S| - \rho|S| + 1$, so two distinct codewords differ on a large
fraction of positions.

The parameter $\rho$ is the \emph{rate}, and $1/\rho$ is the
\emph{blowup factor}. In our implementation the blowup factor is~$2$
($\rho = 1/2$): the codeword produced by the low-degree extension is
twice the length of the coefficient vector.


\subsection{Relative Hamming Distance}
\label{subsec:hamming}

% BS+18 Section 2.1: "the case that f is 0.1-far from (all members of)
% RS[F, S, rho] in relative Hamming distance."
\begin{definition}[Relative Hamming Distance]
\label{def:hamming}
The \emph{relative Hamming distance} between two vectors
$u, v \in \F^n$ is
%
\[
\Delta(u, v) = \frac{|\{i : u_i \neq v_i\}|}{n}.
\]
%
The \emph{distance} of a vector $u$ from a code $C$ is
$\Delta(u, C) = \min_{v \in C} \Delta(u, v)$. A vector $u$ is
$\delta$-\emph{close} to $C$ if $\Delta(u, C) \leq \delta$, and
$\delta$-\emph{far} if $\Delta(u, C) > \delta$.
\end{definition}

\begin{example}
\label{ex:hamming}
The codeword $u = (5,1,4,0,3,6)$ of \cref{ex:rs_code} and the corrupted
word $v = (5,1,4,2,3,6)$ differ in exactly one of six positions, so
$\Delta(u, v) = 1/6$. If every genuine codeword differs from $v$ in at
least this many positions, then $v$ is $\tfrac16$-far from the code.
\end{example}

The FRI protocol tests proximity to a Reed-Solomon codeword. As stated
in~\cite[Section~2.1]{BBHR18}, the verifier's task is to distinguish
between the case that $f \in \RS[\F, S, \rho]$ and the case that $f$
is far from $\RS[\F, S, \rho]$ in relative Hamming distance.


% ==================================================================
\section{Cryptographic Hash Functions}
\label{sec:hash_functions}
% ==================================================================

\begin{definition}[Collision-Resistant Hash Function]
\label{def:hash}
A \emph{collision-resistant hash function} is a function
$H : \{0,1\}^* \to \{0,1\}^\lambda$ such that no efficient algorithm
can find distinct inputs $x \neq x'$ with $H(x) = H(x')$, except
with negligible probability.
\end{definition}

\begin{example}
\label{ex:hash}
SHA-256 maps any input to a $256$-bit digest. Changing even one bit of
the input produces an essentially unrelated digest, and no method is
known to find two inputs with the same digest faster than the roughly
$2^{128}$ tries of a birthday attack, which is what collision resistance
requires.
\end{example}

Hash functions serve two roles in STARK systems: constructing Merkle
trees for committing to polynomial evaluations, and deriving the
verifier's challenges via the Fiat-Shamir heuristic. Our
implementation uses two different hash functions for these two roles:
SHA-256 for the Merkle commitments and Keccak-$f$[1600] for the
challenger.

% CO25 proves DSFS (Duplex Sponge Fiat-Shamir) preserves soundness,
% knowledge-soundness and zero-knowledge.
The distinction matters for the Fiat-Shamir transformation. Chiesa and
Orr\`{u}~\cite{CO25} prove that a \emph{duplex sponge} construction
(such as Keccak/SHA-3) is suitable for Fiat-Shamir, while
Merkle-Damg\r{a}rd constructions (such as SHA-256) lack a comparable
security proof in this setting. For this reason we use a Keccak-based
sponge for the challenger, while retaining SHA-256 for the Merkle
trees, where only collision resistance is required.


% ==================================================================
\section{Proof Systems}
\label{sec:proofs}
% ==================================================================

We now establish the formal framework of proof systems, following
Thaler~\cite{Thaler22}.


\subsection{Interactive Proofs}
\label{subsec:interactive_proofs}

% Thaler, Definition 3.1.
\begin{definition}[Interactive Proof {\cite[Definition~3.1]{Thaler22}}]
\label{def:interactive_proof}
An \emph{interactive proof system} for a language $\mathcal{L}$ is a
pair $(\Prover, \Verifier)$ where $\Prover$ is a computationally
unbounded prover and $\Verifier$ is a probabilistic polynomial-time
verifier. The protocol proceeds in rounds, and must satisfy:
\begin{itemize}
  \item \textbf{Completeness.} For every $x \in \mathcal{L}$,
    $\Pr[\Verifier \text{ accepts}] \geq 1 - \delta_c$.
  \item \textbf{Soundness.} For every $x \notin \mathcal{L}$ and every
    cheating prover $\Prover^*$,
    $\Pr[\Verifier \text{ accepts}] \leq \delta_s$.
\end{itemize}
\end{definition}

\begin{example}
\label{ex:interactive_proof}
A classic interactive proof is for graph non-isomorphism. The verifier
secretly picks one of two graphs, randomly relabels its vertices, and
sends the result to the prover, who must say which graph it came from.
If the graphs are non-isomorphic an unbounded prover always answers
correctly; if they are isomorphic no prover can do better than
guessing, so the verifier catches a false claim with probability
$1/2$ per round.
\end{example}


\subsection{Interactive Arguments}
\label{subsec:arguments}

Relaxing the soundness of an interactive proof to hold only against
computationally efficient provers yields a weaker but more flexible
notion.

\begin{definition}[Interactive Argument {\cite[Section~3.2]{Thaler22}}]
\label{def:argument}
An \emph{interactive argument} for a language $\mathcal{L}$ is an
interactive proof in which the soundness condition is only required to
hold against prover strategies that run in probabilistic polynomial
time.
\end{definition}

\begin{example}
\label{ex:argument}
The FRI-based systems in this thesis are arguments, not proofs. Their
soundness rests on the collision resistance of a hash function: a
prover with unlimited time could find hash collisions and forge a
convincing transcript, but a realistic polynomial-time prover cannot,
so soundness holds against every efficient prover.
\end{example}

Both models share the same structure. In a \emph{public-coin} protocol,
every verifier message is a uniformly random challenge, independent of
the prover's messages, and the verifier's decision is a deterministic
function of the \emph{transcript}, the ordered sequence of messages
and challenges. These are exactly the protocols to which the
Fiat-Shamir transformation (\cref{subsec:fiat_shamir}) applies. An
\emph{argument of knowledge} further requires an efficient extractor
$\Extractor$ that recovers a valid witness from any convincing prover.

STARKs are \emph{arguments}, not \emph{proofs} in the formal sense:
their soundness relies on the collision resistance of a hash function,
which a computationally unbounded prover could break. The interactive
oracle proofs of \cref{sec:iop} specialise this model by providing the
prover's messages as oracles.

\subsection{The Fiat-Shamir Transformation}
\label{subsec:fiat_shamir}

A \emph{non-interactive} proof consists of a single message from the
prover to the verifier. The Fiat-Shamir transformation removes
interaction from a public-coin protocol by replacing the verifier's
random challenges with the output of a hash function applied to the
transcript so far. Under the random oracle model, this preserves
soundness. Chiesa and Orr\`{u}~\cite{CO25} give a formal treatment
using a duplex sponge, which resolves the ambiguity between writing to
and reading from the transcript by using the sponge's native
\emph{absorb} and \emph{squeeze} operations. Our challenger implements
this construction over Keccak-$f$[1600].


% ==================================================================
\section{Polynomial Commitment Schemes}
\label{sec:pcs}
% ==================================================================


\begin{definition}[Polynomial Commitment Scheme]
\label{def:pcs}
A \emph{Polynomial Commitment Scheme} (PCS) consists of four
algorithms:
\begin{itemize}
  \item $\Setup(\lambda, d) \to \mathsf{pp}$: generates public
    parameters for polynomials of degree at most $d$.
  \item $\Commit(\mathsf{pp}, f) \to C$: produces a commitment $C$
    to the polynomial $f$.
  \item $\Open(\mathsf{pp}, f, z) \to (v, \pi)$: produces a proof
    $\pi$ that $f(z) = v$.
  \item $\Verify(\mathsf{pp}, C, z, v, \pi) \to \{0,1\}$: checks
    the evaluation proof.
\end{itemize}
The scheme must satisfy correctness and binding.
\end{definition}

 
\begin{example}
\label{ex:pcs}
Let $f(X) = 5 + 2X + 3X^2$ over $\F_{17}$. The prover first commits to
$f$, producing a commitment $C$ that fixes the polynomial without
revealing its coefficients. Later the verifier names a point, say
$z = 4$, and asks for the value of $f$ there. The prover replies with
$v = f(4) = 5 + 2\cdot 4 + 3\cdot 16 = 61 = 10 \pmod{17}$, together with
a proof $\pi$.
 
The proof $\pi$ convinces the verifier that $v = 10$ is the value at
$z = 4$ of the exact polynomial committed in $C$, and not of some other
polynomial the prover might substitute; the verifier becomes convinced
of the single equality $f(4) = 10$ without ever seeing the coefficients
of $f$. The binding property is what makes this meaningful: the prover
cannot open the same commitment $C$ to two different values at the same
point, so it could not also convince the verifier that $f(4) = 11$. How
$\pi$ is actually built and checked, using the FRI protocol, is the
subject of \cref{ch:pcs}.
\end{example}

The STARK system implemented in this thesis uses a \emph{FRI-based
PCS}~\cite{BBHR18}: a polynomial is committed via a Merkle tree over
its low-degree extension, and evaluation proofs are constructed
through the FRI low-degree testing protocol. This approach requires no
trusted setup (transparent) and relies only on hash functions
(post-quantum secure). The FRI protocol is described in detail in
\cref{ch:fri}.